\section{Introduction}\label{sec:introduction}

Cryptocurrency markets absorb information at intraday frequencies. A
regulatory action against an exchange, a macroeconomic announcement,
or a disruption to trading infrastructure can coincide with a sharp
change in volatility, while routine commentary may carry little new
information. For short-horizon forecasting, the operational question
is how to convert a stream of irregularly arriving text into a useful
component of a volatility forecast. This requires attention both to
what an article says and to how its relevance changes after publication.

Two research traditions motivate our approach. The first, rooted in
the point-process literature, represents events as discrete impulses
whose effects accumulate and dissipate over time. Hawkes processes
\citep{HAWKES1971} and their applications in finance \citep{Bacry2015}
provide a natural vocabulary for these dynamics. Neural extensions
allow richer dependence on event history \citep{Mei2017,Zuo2020}.
The second tradition extracts information from financial text using
sentiment classifiers and contextual representations, including FinBERT
\citep{Araci2019}, CryptoBERT \citep{Kulakowski_2023}, and financial
word embeddings \citep{Rahimikia2023}. Such representations complement
the evidence that sentiment and attention can contribute to volatility
forecasting \citep{Audrino2020}.

These perspectives suggest that text representation and temporal
aggregation should be considered together. A fixed-window average
does not distinguish an article published moments ago from one near
the beginning of the window. A common discount factor distinguishes
their ages but assigns the same persistence to every article. Two
articles published in the same hour, one describing an enforcement
action and another offering routine market commentary, may warrant
different magnitudes and durations of forecast adjustment. We ask
whether these two quantities can be learned jointly, at the article
level, from the content of the news.

We propose News Event Decay (NED), a neural architecture that maps
each article's pretrained CryptoBERT embedding into a positive
magnitude $s_k$ and a decay rate $\lambda_k$. The news branch sums
the signals $s_k\exp(-\lambda_k\Delta t_k)$ over a fixed lookback
window and maps their aggregate into a bounded correction to a
time-series forecast. The forecast is multiplicative in realized
variance and additive in its logarithm. Although the aggregation is
inspired by event-decay kernels, NED is a conditional forecasting
model: it does not estimate an event-arrival intensity or a
self-exciting point process.

The separation between the time-series baseline and the news branch
serves two purposes. It asks news to account for forecast information
beyond a baseline already trained on price and volume history, and it
retains an explicit description of how each article enters the
aggregate signal. The final forecast also depends on a fitted residual
loading, an intercept, and clipping. Consequently, article magnitudes
are interpretable within the fitted correction, rather than as
standalone estimates of causal volatility shocks.

We evaluate one-hour and four-hour realized-variance forecasts on
15-minute Bitcoin and Ethereum bars. Six expanding-window test periods
cover January 2023 through December 2025. The price histories used for
training begin in 2017 for Bitcoin and 2018 for Ethereum. CoinDesk
provides the news corpus; a deterministic tag and title filter removes
common price recaps before articles are aligned to forecast origins.
We compare NED with HAR, EGARCH, an EGARCH+NED hybrid, and a TS-only
model that is exactly the first-phase time-series component of NED.

The main results establish a forecast benefit for Bitcoin. Relative
to TS-only, NED reduces QLIKE by 7.38\% at one hour and 10.83\% at
four hours, with corresponding log-MSE reductions of 0.99\% and
3.11\%. The pooled loss differences are significant at the 1\%
level under the Diebold--Mariano test with the small-sample correction
of \citet{Harvey_1997}. Ethereum gains are smaller: QLIKE declines
by 1.32\% and 2.04\%, while log-MSE differences are not significant.
The contrast between assets limits the scope of a general claim
about the forecasting value of this news source.

A second finding concerns where the Bitcoin gains arise. Sorting
four-hour forecasts by the subsequent return scaled by TS-only
volatility reveals larger improvements in both return tails. In
the central deciles, QLIKE improves modestly while log-MSE can
deteriorate. Similarly, adding the news branch to EGARCH improves
log-MSE but worsens QLIKE. News augmentation therefore changes the
distribution of forecast errors in ways that a single accuracy
measure does not fully describe.

A third set of results examines what the article parameters represent.
Across the BTC four-hour test sample, the median estimated half-life
is 42.7 minutes, and larger internal signals tend to decay more
slowly. The case study, event-category comparison, and sentiment
analysis describe this heterogeneity from complementary perspectives.
Macro news receives the largest mean magnitude and longest median
half-life among the twelve categories.
Negative-classified articles receive larger average magnitudes than
positive-classified articles, but sentiment ambiguity does not produce
a stable separation. These are properties of the learned predictive
representation rather than identified structural responses.

Finally, matched experiments examine the design choices behind the
news branch. A constant kernel already improves QLIKE, and learned
decay provides additional gains, with linear decay performing best
in the kernel comparison. Sentiment with learned decay nearly
matches the full-dimensional embedding, while a 256-component PCA
representation improves on it in this run. These comparisons,
together with lookback sensitivity, distinguish the usefulness of
the news mechanism from the optimality of the main specification.
All variants share the same fitted time-series baseline within each
expanding window.

The remainder of the paper proceeds as follows.
Section~\ref{sec:methodology} describes the data, architecture,
training procedure, baselines, and evaluation protocol.
Section~\ref{sec:results} presents the main results and the Z-score
analysis. Section~\ref{sec:economic} studies article persistence, news
categories, and sentiment. Section~\ref{sec:robustness} compares
alternative decay functions and news representations, with lookback
sensitivity reported in the Appendix.
Section~\ref{sec:conclusion} concludes.
